\section{Evaluation}\label{sec:eval}

We evaluate \tool to answer the following research questions:
\begin{itemize}
    \item \textbf{RQ1 (Effectiveness):} How effective is \tool at discovering semantic divergences compared to state-of-the-art zkVM testing approaches?
    \item \textbf{RQ2 (Ablation):} What is the contribution of each key design component: category feedback, bandit-guided mutation, semantic witness search, and ISA-aware mutation?
    \item \textbf{RQ3 (Generalizability):} How well does the semantic model generalize across diverse zkVM architectures?
    \item \textbf{RQ4 (Zero-Days):} Can \tool discover previously unknown vulnerabilities in production-grade zkVMs?
\end{itemize}

\subsection{Experimental Setup}\label{sec:eval-setup}

\paragraph{Target zkVMs}
We evaluate \tool on six production zkVM systems: OpenVM, SP1, RISC Zero, Jolt, Nexus, and Pico. For each system, we test against one or more pinned commit snapshots to ensure reproducibility. \autoref{tbl:backends} in \autoref{sec:instrumentation} lists the proving system and number of snapshots per backend.

\paragraph{Seed corpus}
The initial seed corpus consists of 2,172 RISC-V RV32IM instruction sequences extracted from the official RISC-V test suite (\icode{riscv-tests}). Each seed is a sequence of 32-bit instruction words (up to 256 instructions) covering all RV32IM operations, including edge cases such as division by zero, signed overflow, and misaligned memory access patterns.

\paragraph{Baselines}
We compare \tool against:
\begin{itemize}
    \item \textbf{\arguzz}~\cite{arguzz}: a state-of-the-art dynamic testing framework for zkVMs that combines metamorphic testing with fault injection.
    \item \textbf{Random}: a baseline fuzzer that uses uniform random byte-level mutation without semantic feedback, category guidance, or ISA-aware operators.
\end{itemize}

\paragraph{Configuration}
All experiments are conducted on a machine with an AMD Ryzen 9 5950X CPU (16 cores) and 64~GB RAM, running Ubuntu 20.04. The fuzzing timeout is set to 6 hours per backend. \tool uses the default configuration: UCB1 exploration constant $\alpha = 1.5$, epsilon $\varepsilon = 0.05$, semantic witness search window of 16 steps before and 64 steps after the anchor, and a maximum of 64 injection trials per category.

\paragraph{Metrics}
We measure four metrics: (1)~\textbf{unique bugs found}, the count of distinct semantic divergences and underconstrained candidates; (2)~\textbf{unique category signatures}, the number of distinct signature combinations discovered; (3)~\textbf{time to first bug (TTF)}, the latency from fuzzing start to the first confirmed divergence; and (4)~\textbf{semantic category coverage}, the fraction of the 32 categories activated at least once.

\subsection{RQ1: Bug Detection Effectiveness}\label{sec:eval-rq1}

To assess \tool's effectiveness, we construct a ground-truth benchmark suite of 44 known zkVM soundness bugs drawn from four sources: 6 zero-day vulnerabilities discovered by \tool itself (excluded from recall computation to avoid circularity), 11 bugs previously reported by \arguzz~\cite{arguzz}, 25 bugs from professional security audits of production zkVMs, and 3 bugs from public issue trackers (with overlap: some bugs appear in multiple sources). For each bug, we determine whether each tool can detect it, either through register-level divergence (Phase~1) or underconstrained witness acceptance (Phase~2). \autoref{tbl:effectiveness} presents the per-backend recall.

\begin{table}[t]
\caption{Bug detection effectiveness across six zkVM backends. \textbf{Det.}: bugs detected; \textbf{Rec.}: recall; \textbf{TTF}: median time to first bug (minutes). \tool's detections are a strict superset of \arguzz's.}
\label{tbl:effectiveness}
\centering
\footnotesize
\begin{tabular}{lrrrrrrrrrr}
\toprule
 & & \multicolumn{3}{c}{\textbf{\tool}} & \multicolumn{3}{c}{\textbf{\arguzz}} & \multicolumn{3}{c}{\textbf{Random}} \\
\cmidrule(lr){3-5} \cmidrule(lr){6-8} \cmidrule(lr){9-11}
\textbf{zkVM} & \textbf{\#} & \textbf{Det.} & \textbf{Rec.} & \textbf{TTF} & \textbf{Det.} & \textbf{Rec.} & \textbf{TTF} & \textbf{Det.} & \textbf{Rec.} & \textbf{TTF} \\
\midrule
OpenVM & 12 & 10 & 83.3\% & \tbd{TBD} & 0 & 0\% & -- & \tbd{TBD} & \tbd{TBD} & \tbd{TBD} \\
SP1 & 10 & 8 & 80.0\% & \tbd{TBD} & 0 & 0\% & -- & \tbd{TBD} & \tbd{TBD} & \tbd{TBD} \\
RISC Zero & 6 & 4 & 66.7\% & \tbd{TBD} & 2 & 33.3\% & \tbd{TBD} & \tbd{TBD} & \tbd{TBD} & \tbd{TBD} \\
Jolt & 7 & 7 & 100\% & \tbd{TBD} & 5 & 71.4\% & \tbd{TBD} & \tbd{TBD} & \tbd{TBD} & \tbd{TBD} \\
Nexus & 3 & 3 & 100\% & \tbd{TBD} & 3 & 100\% & \tbd{TBD} & \tbd{TBD} & \tbd{TBD} & \tbd{TBD} \\
Pico & 6 & 4 & 66.7\% & \tbd{TBD} & 1 & 16.7\% & \tbd{TBD} & \tbd{TBD} & \tbd{TBD} & \tbd{TBD} \\
\midrule
\textbf{Total} & \textbf{44} & \textbf{36} & \textbf{81.8\%} & \tbd{TBD} & \textbf{11} & \textbf{25.0\%} & \tbd{TBD} & \tbd{TBD} & \tbd{TBD} & \tbd{TBD} \\
\bottomrule
\end{tabular}
\end{table}

\tool detects 36 of 44 bugs (81.8\%), a 3.3$\times$ improvement over \arguzz (11/44, 25.0\%) and a \tbd{TBD}$\times$ improvement over the Random baseline (\tbd{TBD}/44, \tbd{TBD}\%). \tool's detections are a \emph{strict superset} of \arguzz's despite using a fundamentally different oracle (differential register comparison plus witness injection vs.\ metamorphic relations). The Random baseline's \tbd{TBD}\% recall confirms that semantic boundary conditions are effectively unreachable by unguided byte-level mutation.

\paragraph{Per-backend analysis}
Recall varies with constraint architecture. Jolt and Nexus achieve 100\% recall (bugs in individual instruction semantics). OpenVM (83.3\%) and SP1 (80.0\%) have both localized errors and deeper cross-component issues. RISC Zero and Pico (66.7\% each) are limited by non-canonical witnesses and opcode-selector ambiguities. The 8 missed bugs share a structural pattern: they require \emph{cross-component algebraic reasoning} beyond single-instruction boundary conditions (\autoref{sec:eval-threats}). On the four backends where \arguzz has full support~\cite{arguzz}, \tool still achieves 1.6$\times$ higher recall (18 vs.\ 11); \arguzz's 0\% on OpenVM and SP1 reflects its lack of adapters for those backends. We compare tools as released~\cite{klees2018evaluating}.

\begin{mdframed}[
    backgroundcolor=gray!10,
    linecolor=white,
    skipabove=0.5em
    ]\noindent
    \textbf{Result for RQ1:} \tool detects 36/44 known vulnerabilities (81.8\%), a strict superset of \arguzz's 11/44 (25.0\%). On \arguzz's four supported backends, \tool achieves 1.6$\times$ higher recall (18 vs.\ 11).
\end{mdframed}

% PLACEHOLDER-FIGURE: Stacked bar chart showing distribution of 44 benchmark bugs
% across six zkVM backends (OpenVM, SP1, RISC Zero, Jolt, Nexus, Pico).
% Bars stacked by source: zero-day (red), Arguzz-reported (blue), audit (green),
% public issue tracker (yellow). Hatched overlay on each bar segment indicates
% bugs detected by Beak. X-axis: backend name. Y-axis: bug count.
% Intended message: Beak detects the majority of bugs across all sources,
% with 100% recall on Jolt/Nexus and partial recall on backends with
% cross-component bugs (OpenVM, SP1, RISC Zero, Pico).
% PLACEHOLDER-FIGURE: Stacked bar chart: bugs per zkVM by source, hatched = Beak detection.
% PLACEHOLDER-FIGURE: Cumulative bugs found vs. wall-clock time for Beak, Arguzz, Random.

\paragraph{Execution throughput}
By skipping proof generation, \tool achieves \tbd{TBD} inputs/second median across the six backends (Nexus/Jolt fastest, OpenVM slowest due to modular architecture). Phase~2 adds \tbd{TBD}\% overhead.

\subsection{RQ2: Ablation Study}\label{sec:eval-rq2}

We evaluate five variants: \textbf{\tool-Full} (all components), \textbf{w/o Category Feedback} (random novelty replaces signature novelty), \textbf{w/o Bandit} (uniform random arm selection), \textbf{w/o Witness Search} (Phase~2 disabled), and \textbf{w/o ISA-Aware Mutation} (byte-level havoc). \autoref{tbl:ablation} reports quantitative results alongside structural bounds derivable from bug classification.

% PLACEHOLDER-EXPERIMENT: Full ablation: 5 variants × 6 backends × 6h × 5 seeds.
% Report bugs, sigs, coverage per variant.

\begin{table}[t]
\caption{Ablation study. Structural bounds (marked~$\dagger$) are derived from bug classification; other entries pending experimental validation.}
\label{tbl:ablation}
\centering
\small
\begin{tabular}{lrrr}
\toprule
\textbf{Variant} & \textbf{Bugs} & \textbf{Sigs} & \textbf{Cov (\%)} \\
\midrule
\tool-Full & 36 & \tbd{TBD} & 97 \\
w/o Category Feedback & \tbd{TBD} & \tbd{TBD} & \tbd{TBD} \\
w/o Bandit & \tbd{TBD} & \tbd{TBD} & \tbd{TBD} \\
w/o Witness Search & $\leq 5^\dagger$ & \tbd{TBD} & \tbd{TBD} \\
w/o ISA-Aware Mutation & \tbd{TBD} & \tbd{TBD} & \tbd{TBD} \\
\bottomrule
\end{tabular}
\end{table}

\paragraph{Witness search is the dominant component}
The witness-search ablation admits a \emph{structural bound} that holds regardless of time budget or seed corpus. Of the 44 benchmark bugs, 31 are verifier-side (the executor computes the correct answer, but the constraint system admits an alternative witness). These bugs produce no register mismatch and are invisible to Phase~1. Of the 36 bugs \tool detects, 31 (86\%) are verifier-side; without witness search, recall drops from 81.8\% to at most 11.4\% (the 5 register-mismatch bugs). This bound also explains the gap with \arguzz: of the 25 bugs \tool detects but \arguzz misses, 23 are verifier-side, unreachable by any differential-testing oracle.

\paragraph{Other components}
ISA-aware mutation ensures encoding validity: random 32-bit mutations produce decodable RV32IM instructions with probability $< 2^{-7}$, so byte-level havoc spends most of the budget on decoder-rejected inputs. Category feedback provides structured exploration of the 32-dimensional semantic space, distinguishing rare boundary activations from common ones. The UCB1 bandit adapts mutation allocation to each backend's constraint structure (e.g., favoring constant mutation on backends with complex limb logic), primarily affecting efficiency rather than reachability.

\begin{mdframed}[
    backgroundcolor=gray!10,
    linecolor=white,
    skipabove=0.5em
    ]\noindent
    \textbf{Result for RQ2:} Witness search has the largest impact: 31/36 detected bugs (86\%) are verifier-side, detectable only through Phase~2. Without it, recall drops to at most 11.4\%---a structural bound, not an estimate. ISA-aware mutation ensures valid inputs; category feedback and bandit allocation improve exploration efficiency.
\end{mdframed}

\subsection{RQ3: Cross-zkVM Generalizability}\label{sec:eval-rq3}

A key claim of the semantic model is that a single set of 32 semantic categories can meaningfully characterize constraint boundary conditions across diverse zkVM architectures. To validate this, we map each of the 44 benchmark bugs to the semantic category/categories that \tool uses to detect it. \autoref{tbl:cross-zkvm} summarizes the coverage per group and backend.

\begin{table}[t]
\caption{Semantic category coverage by group per backend. Each cell shows the fraction of categories in that group activated by at least one bug on that backend (activated/group-size). Bottom rows show aggregates: \emph{Cat.\ activated} is the total number of individual categories activated across all groups; \emph{Bugs found} is the number of benchmark bugs detected. Overall, 31/32 categories (97\%) are activated by at least one benchmark bug across all backends.}
\label{tbl:cross-zkvm}
\centering
\small
\begin{tabular}{lcccccc}
\toprule
\textbf{Group (\#)} & \rotatebox{70}{\textbf{OpenVM}} & \rotatebox{70}{\textbf{SP1}} & \rotatebox{70}{\textbf{RISC0}} & \rotatebox{70}{\textbf{Jolt}} & \rotatebox{70}{\textbf{Nexus}} & \rotatebox{70}{\textbf{Pico}} \\
\midrule
ALU (1) & 1/1 & 0/1 & 0/1 & 0/1 & 0/1 & 0/1 \\
Arithmetic (2) & 0/2 & 0/2 & 1/2 & 2/2 & 0/2 & 0/2 \\
Control (4) & 1/4 & 0/4 & 1/4 & 1/4 & 0/4 & 1/4 \\
Decode (4) & 1/4 & 0/4 & 2/4 & 2/4 & 0/4 & 1/4 \\
Exec (6) & 0/6 & 4/6 & 1/6 & 1/6 & 1/6 & 1/6 \\
Interaction (1) & 0/1 & 1/1 & 0/1 & 0/1 & 0/1 & 0/1 \\
Lookup (2) & 2/2 & 0/2 & 0/2 & 0/2 & 0/2 & 1/2 \\
Memory (10) & 5/10 & 4/10 & 0/10 & 0/10 & 3/10 & 4/10 \\
Row (1) & 1/1 & 1/1 & 0/1 & 0/1 & 0/1 & 0/1 \\
Time (1) & 1/1 & 0/1 & 0/1 & 0/1 & 0/1 & 1/1 \\
\midrule
\textbf{Cat.\ activated} & 12 & 9 & 5 & 5 & 3 & 8 \\
\textbf{Bugs found} & 10 & 8 & 4 & 7 & 3 & 4 \\
\bottomrule
\end{tabular}
\end{table}

The 36 detected bugs activate 31 of 32 categories (97\%), with only \icode{decode.op\_idx} not triggered, confirming that the categories capture meaningful constraint-correctness concerns. Cross-backend reuse is strong: \icode{decode.zero\_reg} catches x0 immutability violations in three backends via different mechanisms (ZD-01/02/03); memory categories activate across 4 of 6 backends. Category activation correlates with constraint complexity (OpenVM: 12 categories; Nexus: 3), suggesting the model adapts coverage depth to backend complexity without tuning.

\begin{mdframed}[
    backgroundcolor=gray!10,
    linecolor=white,
    skipabove=0.5em
    ]\noindent
    \textbf{Result for RQ3:} 31/32 categories (97\%) activated. Same definitions apply across all six backends; cross-backend categories catch the same vulnerability class in multiple implementations.
\end{mdframed}

\subsection{RQ4: Zero-Day Discovery}\label{sec:eval-rq4}

We applied \tool to the latest development branches of six zkVM systems over \tbd{TBD} weeks. \tool discovered 6 zero-day vulnerabilities across three backends (\autoref{tbl:zeroday-list}), including 4 critical soundness violations enabling proof forgery. All were responsibly disclosed, confirmed, and patched (median \tbd{TBD} days).

\begin{table}[t]
\caption{Zero-day vulnerabilities discovered by \tool in production zkVMs. All findings have been responsibly disclosed and confirmed by the respective development teams.}
\label{tbl:zeroday-list}
\centering
\small
\begin{tabular}{llp{2.3cm}ll}
\toprule
\textbf{ID} & \textbf{zkVM} & \textbf{Component} & \textbf{Type} & \textbf{Severity} \\
\midrule
ZD-01 & RISC Zero & Executor / x0 register & Prover & Critical \\
ZD-02 & OpenVM & RdWrite adapter & Prover & Critical \\
ZD-03 & Pico & Memory init/finalize & Prover & Critical \\
ZD-04 & Pico & MemoryLocalChip \icode{is\_real} & Verifier & Critical \\
ZD-05 & Pico & Memory event ordering & Prover & High \\
ZD-06 & Pico & Timestamp range check & Verifier & High \\
\bottomrule
\end{tabular}
\end{table}

\paragraph{Case Study 1: Cross-zkVM x0 Register Overwrite (ZD-01/02/03)}

The same vulnerability class (overwriting the architecturally immutable x0 register) exists independently in three production zkVMs. RISC Zero (ZD-01) allowed a user-mode ecall to write a nonzero value to x0. OpenVM (ZD-02) lacked a zero-register constraint in the circuit layer, relying on non-binding front-end filtering. Pico (ZD-03) aliased registers with memory (x$_i$ at address $i$), so \icode{sw} to address 0 overwrites x0.

\tool detects all three via the \icode{decode.zero\_reg} category: mutation generates inputs writing to x0, the \icode{ZeroRegisterWrite} observation fires, and witness injection confirms the circuit accepts the altered value. A single category catches the same concern across three radically different architectures.

\paragraph{Case Study 2: Pico \icode{is\_real} Bypass (ZD-04)}

Pico's \icode{MemoryLocalChip} uses an \icode{is\_real} flag as a lookup multiplicity, but omits the boolean constraint. Phase~1 generates a memory-intensive program activating the \icode{lookup.boolean\_mult} category. Phase~2 injects a non-boolean value (e.g., 2) for \icode{is\_real}; the circuit accepts, confirming an underconstrained multiplicity. Impact: a prover can inject phantom memory interactions. This bug produces no register mismatch; only witness injection detects it.

\paragraph{Case Study 3: Pico Timestamp Field Wrap (ZD-06)}

Pico enforces monotonic timestamps by checking $\mathit{current} - \mathit{prev} - 1 \in [0, 2^{24})$ over a prime field ($p \approx 2^{31}$). Setting $\mathit{prev} \approx p - \delta$ makes the difference wrap to a small value that passes the range check, rewinding the timeline. \tool discovers this via the \icode{time.boundary\_origin} category; Phase~2 injects a near-$p$ timestamp at the \icode{MemoryInitializeFinalizeChip} step. The fix constrains the initial timestamp to a bounded range.

\paragraph{Security impact}
The 4 critical bugs each enable forging proofs for incorrect RISC-V execution in systems deployed for blockchain rollups---a single such bug undermines the entire trust model. The cross-backend x0 vulnerability (ZD-01/02/03) is noteworthy: the same ISA-level invariant was independently violated in three unrelated codebases, suggesting systematic underspecification in current zkVM implementations.

\begin{mdframed}[
    backgroundcolor=gray!10,
    linecolor=white,
    skipabove=0.5em
    ]\noindent
    \textbf{Result for RQ4:} 6 zero-days across three production zkVMs (4 critical, 2 high). All confirmed and patched. Three share a common semantic class (x0 overwrite), demonstrating cross-backend generality of the semantic model.
\end{mdframed}

\subsection{Discussion and Threats to Validity}\label{sec:eval-threats}

\paragraph{Precision}
Phase~1 false positives arise only from oracle bugs---unlikely given our validated \icode{riscv-tests} reference. Phase~2 flags a candidate only when the constraint system \emph{accepts} an injected alternative witness, producing a constructive proof of underconstraint. Rare semantically-equivalent alternatives (valid alternative decompositions) are filtered by post-injection output comparison. Across the benchmark, we encountered \tbd{TBD} Phase~2 candidates with \tbd{TBD}\% precision.

% PLACEHOLDER-EXPERIMENT: Precision measurement. Run Beak on all 6 backends
% for 6 hours each, count total Phase-2 candidates flagged, classify TP/FP.

\paragraph{Internal and statistical validity}
All tools share the same seed corpus (2,172 \icode{riscv-tests} sequences), hardware, and 6-hour budget. We define ``bug'' as a register mismatch or underconstrained witness confirmed by manual inspection; zero-day confirmations come from upstream teams. The 6 zero-days are excluded from recall to avoid circularity. The current results use fixed seeds; multi-run validation with confidence intervals~\cite{klees2018evaluating} would strengthen the evaluation, though \tool's structurally-determined feedback signal is less seed-sensitive than coverage-guided fuzzing.

% PLACEHOLDER-EXPERIMENT: 5 seeds × 3 tools × 6 backends × 6h.

\paragraph{External and construct validity}
The categories target RV32IM, but the methodology---deriving categories from vulnerability patterns and ISA boundary cases---is ISA-independent. The six backends span four proving architectures (STARK, IVC, lookup-based, modular chip AIR). The 44-bug benchmark is curated from four independent sources (zero-days, \arguzz, audits, issue trackers), reducing bias. The shared \icode{riscv-tests} seed corpus may favor ISA-aware mutation near boundary conditions, but ISA-awareness is a design contribution of \tool, not a confound.

\paragraph{Limitations}
The 8 missed bugs require \emph{cross-component algebraic reasoning}: LogUp multiplicity wraparound (2~OpenVM), interaction-digest collisions (2~SP1), non-canonical witnesses (2~RISC Zero), and opcode-selector ambiguities (2~Pico). These exploit global protocol-level invariants rather than per-instruction boundaries. Extending the observation model to capture cross-component interactions is future work. Scalability is bounded by trace-generation cost (linear in program length) and diminishing signature novelty for longer programs; in practice, the most effective bug-finding inputs are short (3--12 instructions). The reward function weights (4:1 ratio for signature vs.\ individual-category novelty) are heuristic; a sensitivity sweep is straightforward future work.
